% !TeX root=../main.tex

\chapter{مقدمه}
\label{chap:introduction}

% ------------------------------------------------------------
\section{انگیزه و بیان مسئله}
\label{sec:motivation}
% ------------------------------------------------------------

پروژه‌های بزرگ زیرساختی و مهندسی معمولاً از بودجه و زمان‌بندی
پیش‌بینی‌شده فراتر می‌روند \citep{Flyvbjerg2003, Merrow2011}.
وقتی چند پروژه از این نوع هم‌زمان و با یک سرمایه‌ی مشترک اجرا شوند،
ریسک‌ها روی هم انباشته می‌شوند. ضعف عملکرد در یک پروژه از راه
بودجه‌ی مشترک به پروژه‌های دیگر هم منتقل می‌شود، و اینکه در هر دوره
چه مقدار پول به هر پروژه برسد می‌تواند تعیین کند که سبد پروژه‌ها
به سرانجام برسد یا شکست بخورد.

برای درک عمق این مسئله کافی است به محیط واقعی اجرای پروژه‌های بزرگ
نگاه کنیم. سازمان‌هایی که به‌طور هم‌زمان چند پروژه‌ی بزرگ را اجرا
می‌کنند، معمولاً با یک منبع مالی مرکزی کار می‌کنند که باید میان
پروژه‌های مختلف تقسیم شود. این منبع می‌تواند خط اعتباری بانکی،
بودجه‌ی سالانه‌ی سازمانی، یا سرمایه‌ی اختصاص‌یافته از سوی سرمایه‌گذار
باشد. در هر صورت، محدودیت مشترک است و پروژه‌ها برای دریافت سهم
بیشتری از آن با یکدیگر رقابت می‌کنند.

این رقابت در شرایط عادی قابل مدیریت است. مشکل زمانی بروز می‌کند که
یک یا چند پروژه دچار تأخیر یا افزایش هزینه می‌شوند. مدیر سبد در
چنین شرایطی با یک معضل واقعی روبه‌روست: آیا باید بودجه‌ی بیشتری به
پروژه‌ی عقب‌افتاده تخصیص داد تا آن را نجات داد، یا باید منابع را به
پروژه‌های سالم‌تر هدایت کرد تا از انتقال بیماری جلوگیری شود؟ هیچ
پاسخ ساده‌ای برای این پرسش وجود ندارد، چون هر انتخاب پیامدهای زنجیره‌ای
در سراسر سبد دارد.

پیچیدگی بیشتر می‌شود وقتی می‌دانیم که بازده‌ی بودجه در پروژه‌های
عمرانی و مهندسی اساساً نامطمئن است. یک مقدار مشخص پول در یک دوره
ممکن است در شرایط مساعد پیشرفت زیادی ایجاد کند، و در همان پروژه در
دوره‌ی بعد، به دلیل شرایط آب‌وهوایی، مشکلات تأمین مواد، یا ضعف
عملکرد نیروی کار، پیشرفت بسیار کمتری داشته باشد. این عدم‌قطعیت
ذاتی است و با دقت بیشتر در برنامه‌ریزی از بین نمی‌رود؛ فقط می‌توان
آن را به‌درستی مدل کرد و در تصمیم‌گیری لحاظ نمود.

با این حال، تخصیص بودجه در سطح سبد، به‌عنوان یک مسئله‌ی تصمیم‌گیری
ترتیبی، در ادبیات مدیریت سبد پروژه تقریباً بررسی نشده است.
بیشتر پژوهش‌ها بر مرحله‌ی \textbf{انتخاب} تمرکز کرده‌اند، یعنی اینکه
کدام پروژه‌ها وارد سبد شوند. اما مرحله‌ی \textbf{راهبری عملیاتی}،
یعنی اینکه سبدِ انتخاب‌شده تا پایان کار چگونه اداره شود،
کمتر مورد تحلیل قرار گرفته است
\citep{Archer1999, Cooper2001, Turner2014}.
این پژوهش با ساخت یک محیط شبیه‌سازی صحت‌سنجی‌شده به این کمبود
می‌پردازد. این محیط مسئله را برای روش‌های داده‌محور قابل حل می‌کند
و با ایجاد یک خط مبنای عملکرد نشان می‌دهد تخصیص ساده‌انگارانه
چه هزینه‌ای دارد.

این شکاف در ادبیات تصادفی نیست. ریشه‌ی آن در نحوه‌ی شکل‌گیری رشته‌ی
مدیریت پروژه است. مدیریت پروژه به‌عنوان یک رشته‌ی آکادمیک، عمدتاً
از مسائل برنامه‌ریزی و زمان‌بندی رشد کرده است. روش مسیر بحرانی و
تکنیک ارزیابی و بازنگری برنامه، هر دو ابزارهایی برای مدیریت یک
پروژه‌ی منفرد هستند. انتقال این چارچوب‌ها به مدیریت چند پروژه‌ی
هم‌زمان، آن هم زیر یک قید بودجه‌ی مشترک، چالش‌هایی ایجاد می‌کند که
ابزارهای کلاسیک برای آن‌ها طراحی نشده‌اند.

از سوی دیگر، ادبیات مدیریت مالی و نظریه‌ی سبد نیز، با وجود پیشرفت‌های
قابل توجه در مدل‌سازی ریسک و بهینه‌سازی پرتفولیو، اساساً به دارایی‌های
مالی مربوط می‌شود، نه به پروژه‌های عمرانی با چرخه‌ی اجرایی چندساله.
پروژه‌ی ساختمانی یا زیرساختی ویژگی‌هایی دارد که آن را از سهام و
اوراق قرضه متمایز می‌کند: غیرقابل تقسیم بودن، وابستگی به عملکرد
فیزیکی، قراردادهای پیچیده با طرف مقابل، و هزینه‌های فسخ قابل توجه.
این ویژگی‌ها ایجاب می‌کنند که رویکرد مدل‌سازی متفاوتی اتخاذ شود.

ما سبدی از پروژه‌ها را در نظر می‌گیریم که هر کدام با یک قرارداد
دوطرفه مطابق فرم‌های \lr{FIDIC} \citep{FIDIC2017} اداره می‌شود و همه
با یک ذخیره‌ی سرمایه‌ی مشترک هم‌زمان پیش می‌روند.
مدیر سبد در هر دوره‌ی برنامه‌ریزی گسسته یک تصمیم تکراری می‌گیرد:
چه مقدار بودجه به هر پروژه‌ی فعال بدهد.
مسئله این است که سیاست تخصیصی پیدا شود که جریان نقدی خالص
تنزیل‌شده‌ی سبد را در کل افق برنامه‌ریزی بیشینه کند و در همان حال
قواعد قرارداد را رعایت کند.

انتخاب قراردادهای \lr{FIDIC} به‌عنوان چارچوب قراردادی این پژوهش از
روی دقت و قصد بوده است. فرم‌های \lr{FIDIC} در پروژه‌های زیرساختی
بین‌المللی، به‌خصوص در کشورهای در حال توسعه و پروژه‌های تأمین‌مالی‌شده
توسط نهادهای چندجانبه مثل بانک جهانی و بانک توسعه‌ی آسیایی، به‌طور
گسترده به کار می‌روند. این فرم‌ها ساختار مالی پیچیده‌ای دارند که
شامل پیش‌پرداخت، پرداخت‌های مرحله‌ای، کسر سپرده‌ی تضمین و بندهای فسخ
است. این پیچیدگی دقیقاً چیزی است که مدل‌های ساده‌انگارانه‌ی تخصیص
بودجه نادیده می‌گیرند.

در نهایت، هدف این پژوهش نه ارائه‌ی یک توصیه‌ی مدیریتی کلی، بلکه
ساخت یک بستر محاسباتی دقیق است که بتواند پیچیدگی واقعی مسئله را
بازنمایی کند. بدون چنین بستری، هر ادعایی درباره‌ی بهینه‌بودن یک
سیاست تخصیص فقط یک ادعای کیفی است. با این بستر، می‌توان سیاست‌های
مختلف را به‌صورت کمّی مقایسه کرد و گفت که کدام یک در چه شرایطی
بهتر عمل می‌کند.

% ------------------------------------------------------------
\section{زمینه و اهمیت موضوع}
\label{sec:context}
% ------------------------------------------------------------

برای درک اهمیت مسئله‌ی تخصیص بودجه در سطح سبد، لازم است ابتدا
تصویری روشن از محیطی که این مسئله در آن رخ می‌دهد ارائه شود.
پروژه‌های زیرساختی و مهندسی بزرگ، اعم از ساخت بزرگراه، احداث
نیروگاه، توسعه‌ی شبکه‌ی فاضلاب شهری، یا ساخت بیمارستان، ماهیتاً
پروژه‌هایی هستند که به دهه‌های طولانی برنامه‌ریزی، چندین سال اجرا،
و منابع مالی کلان نیاز دارند. سازمان‌هایی که این پروژه‌ها را مدیریت
می‌کنند، معمولاً شرکت‌های دولتی، کنسرسیوم‌های بین‌المللی، یا نهادهای
توسعه‌ای هستند که باید چندین پروژه را به‌طور هم‌زمان پیش ببرند.

چالش اصلی این سازمان‌ها آن است که بودجه‌ی کل آن‌ها در هر دوره‌ی
زمانی محدود است، اما تعهداتشان به پروژه‌های مختلف پیوسته و قراردادی
است. نمی‌توانند بگویند که این ماه به پروژه‌ی الف پول نمی‌دهند چون
پروژه‌ی ب ضروری‌تر است، چون این تصمیم ممکن است موجب فسخ قرارداد
پروژه‌ی الف و بروز خسارات قراردادی شود. از طرف دیگر، توزیع یکنواخت
بودجه بین همه‌ی پروژه‌ها هم لزوماً بهینه نیست، چون پروژه‌های مختلف
در مراحل مختلفی از اجرا قرار دارند و نیازهای مالی متفاوتی دارند.

یکی از ویژگی‌های مهم این مسئله، ماهیت ترتیبی تصمیم‌گیری است.
تصمیم تخصیص در هر دوره نه‌تنها بر نتیجه‌ی آن دوره بلکه بر وضعیت
پروژه‌ها در دوره‌های آینده هم اثر می‌گذارد. اگر این ماه به پروژه‌ای
کمتر از نیازش بودجه داده شود، ممکن است در ماه‌های آینده هزینه‌ی
جبران‌سازی بیشتری لازم باشد. یا بدتر از آن، ممکن است پروژه به
وضعیتی برسد که دیگر قابل نجات نباشد و فسخ آن اجتناب‌ناپذیر شود.
این وابستگی زمانی تصمیم‌ها، مسئله را از یک مسئله‌ی بهینه‌سازی
ایستا به یک مسئله‌ی تصمیم‌گیری پویا تبدیل می‌کند.

ویژگی دیگر این مسئله، وجود عدم‌قطعیت در سطوح مختلف است. اول،
بازده‌ی بودجه به پیشرفت نامطمئن است؛ همان‌طور که گفته شد، یک مقدار
مشخص پول در شرایط مختلف پیشرفت‌های مختلفی ایجاد می‌کند. دوم،
وضعیت آینده‌ی پروژه‌ها به مجموعه‌ی تصمیم‌های گذشته و تحقق‌های
تصادفی آن دوره‌ها وابسته است. سوم، افق زمانی پروژه‌ها کاملاً
قطعی نیست؛ تأخیرها می‌توانند افق را طولانی‌تر کنند و پیشرفت
بهتر از انتظار می‌تواند آن را کوتاه‌تر کند. این سه لایه از
عدم‌قطعیت با هم ترکیب می‌شوند و محیطی می‌سازند که برای روش‌های
کلاسیک تحقیق در عملیات به‌شدت چالش‌برانگیز است.

در چنین محیطی، تکیه بر قضاوت شهودی مدیر یا قوانین سرانگشتی
ساده، هزینه‌ای قابل اندازه‌گیری دارد. این پژوهش نشان می‌دهد که
این هزینه با افزایش پیچیدگی سبد به‌طور منظم بزرگ‌تر می‌شود.
به عبارت دیگر، هر چه سبد بزرگ‌تر و پیچیده‌تر باشد، شکاف بین
تصمیم‌گیری سرانگشتی و تصمیم‌گیری بهینه عمیق‌تر می‌شود. این
نتیجه، دلیل اصلی نیاز به رویکردهای مبتنی بر یادگیری ماشین برای
این مسئله است.

\subsection{پروژه‌های زیرساختی و ویژگی‌های خاص آن‌ها}
\label{subsec:infra_characteristics}

پروژه‌های زیرساختی به دلایل متعددی با سایر انواع پروژه متفاوتند.
از منظر مدیریت مالی، چند ویژگی آن‌ها برای این پژوهش اهمیت ویژه دارد.

نخست، \textbf{غیرقابل انقسام بودن} این پروژه‌هاست. برخلاف دارایی‌های
مالی که می‌توان کسری از آن‌ها را خرید یا فروخت، یک پروژه‌ی ساختمانی
یا کامل می‌شود یا نمی‌شود. نمی‌توان «نصف» یک بزرگراه را بهره‌برداری
کرد یا «سه‌چهارم» یک نیروگاه را راه‌اندازی نمود. این ویژگی به این
معناست که ارزش پروژه به‌صورت تدریجی تا لحظه‌ی تکمیل آزاد نمی‌شود،
بلکه عمدتاً در لحظه‌ی تحویل و بهره‌برداری محقق می‌شود. در نتیجه،
جریان نقدی برگشتی پروژه ماهیت پله‌ای دارد، نه پیوسته.

دوم، \textbf{ماهیت قراردادی} روابط مالی است. مدیر سبد مستقیماً کار
پروژه را انجام نمی‌دهد؛ او از طریق قرارداد با پیمانکار کار را
واگذار می‌کند. این قرارداد تعهدات مالی مشخصی برای هر دو طرف ایجاد
می‌کند. مدیر سبد موظف است طبق برنامه‌ی پرداخت مندرج در قرارداد،
مبالغ تعیین‌شده را به پیمانکار پرداخت کند. اگر این پرداخت‌ها با
تأخیر صورت گیرند یا ناکافی باشند، پیمانکار از نظر قراردادی
می‌تواند ادعای خسارت مطرح کند یا در موارد شدید، قرارداد را فسخ
کند.

سوم، \textbf{هزینه‌های فسخ} قابل توجه است. برخلاف فروش یک سهام که
با یک تراکنش ساده انجام می‌شود، فسخ یک قرارداد پروژه فرایندی
پیچیده است که هزینه‌های مستقیم و غیرمستقیم قابل توجهی دارد.
هزینه‌های مستقیم شامل پرداخت غرامت به پیمانکار برای کار انجام‌شده‌ای
است که هنوز صورت‌وضعیت نشده، هزینه‌های بسیج و انحلال کارگاه، و
احتمالاً خسارت تأخیر که در قرارداد مشخص شده است. هزینه‌های غیرمستقیم
شامل اختلال در برنامه‌ریزی کلی سازمان، آسیب به شهرت در بازار
پیمانکاری، و از دست رفتن ارزش کارهای نیمه‌تمامی است که ممکن است
قابل استفاده برای پیمانکار جدید نباشند.

چهارم، \textbf{وابستگی به پیشرفت فیزیکی} به‌عنوان معیار پرداخت
است. در قراردادهای مبتنی بر نقطه‌ی عطف مثل آنچه در \lr{FIDIC}
تعریف شده، پرداخت زمانی صورت می‌گیرد که پروژه به آستانه‌های پیشرفت
از پیش تعیین‌شده برسد. این بدان معناست که اگر پیمانکار نتواند به
نقطه‌ی عطف بعدی برسد، پول جدیدی دریافت نمی‌کند. این ساختار
انگیزه‌ی پیمانکار را با هدف مدیر سبد هم‌راستا می‌کند، اما در عین
حال پیش‌بینی زمان دریافت پول را وابسته به عملکرد واقعی پروژه می‌کند
که ذاتاً نامطمئن است.

\subsection{بافت مالی و محدودیت‌های سازمانی}
\label{subsec:financial_context}

سازمان‌هایی که سبدهای بزرگ پروژه را مدیریت می‌کنند معمولاً در یک
بافت مالی خاص عمل می‌کنند که باید در مدل‌سازی لحاظ شود. این بافت
شامل چند عنصر کلیدی است.

\textbf{جریان نقدی دوطرفه:} بودجه‌ی سبد نه یک موجودی ثابت بلکه یک
جریان پویاست. از یک سو، سازمان پول خرج می‌کند: به پیمانکارها
پیش‌پرداخت می‌دهد، صورت‌وضعیت‌های نقطه‌ی عطف را می‌پردازد، و
هزینه‌های اداری سبد را تأمین می‌کند. از سوی دیگر، پول دریافت
می‌کند: از کارفرمای اصلی یا تأمین‌مالی‌کننده، از کارمزد مدیریت
پروژه، و در برخی موارد از محل بهره‌ی موجودی نقدی. سرعت ورود و
خروج پول در هر دوره‌ی زمانی تعیین می‌کند که آیا سبد می‌تواند به
تعهداتش عمل کند یا خیر.

\textbf{هزینه‌ی فرصت سرمایه:} پولی که به یک پروژه تخصیص می‌یابد
دیگر در دسترس پروژه‌ی دیگر نیست. در شرایطی که بازار فرصت‌های
سرمایه‌گذاری دارد، این هزینه‌ی فرصت واقعی است. ضریب تنزیل در مدل
ما نه فقط ترجیح زمانی بلکه همین هزینه‌ی فرصت را هم منعکس می‌کند.
سازمانی که می‌توانست با همان پول یک پروژه‌ی جدید شروع کند اما آن
را در یک پروژه‌ی کند گیر انداخته، یک هزینه‌ی واقعی دارد حتی اگر
این هزینه در ترازنامه ظاهر نشود.

\textbf{قید عدم استقراض:} در بسیاری از قراردادها و مقررات سازمانی،
موجودی نقدی سبد نمی‌تواند منفی شود. این قید به این معناست که
سازمان نمی‌تواند بیش از آنچه در اختیار دارد تخصیص دهد. وجود این
قید، مسئله‌ی تخصیص را از یک مسئله‌ی بهینه‌سازی بدون قید تبدیل به
یک مسئله‌ی برنامه‌ریزی محدودشده می‌کند.

\textbf{ورودی‌های دوره‌ای:} در اغلب پروژه‌های زیرساختی، کارفرما یا
تأمین‌مالی‌کننده بودجه را به‌صورت دوره‌ای و طبق یک برنامه‌ی از پیش
تعیین‌شده منتقل می‌کند. این ورودی‌های دوره‌ای، به‌علاوه‌ی پرداخت‌های
دریافتی از پیمانکاران بابت تسویه‌ی پیش‌پرداخت‌ها، جریان ورودی
نقدی سبد را می‌سازند. سازمان نمی‌تواند این ورودی‌ها را کنترل کند
و باید با آنچه هر دوره دریافت می‌کند بهترین تخصیص را انجام دهد.

% ------------------------------------------------------------
\section{سلسله‌مراتب تصمیم‌گیری در مدیریت پروژه}
\label{sec:hierarchy}
% ------------------------------------------------------------

کار حاضر در لایه‌ی راهبری سبد قرار می‌گیرد که یکی از سه لایه‌ی
سلسله‌مراتب مدیریت پروژه است \citep{PMI2017}.
انتخاب سبد، یعنی تصمیم درباره‌ی اینکه به کدام پروژه‌ها متعهد شویم،
یک ورودی بیرونی فرض می‌شود. این فرض با مدل راهبری دومرحله‌ای
\citep{Turner2014} سازگار است.
در این لایه مدیر سبد اجرای پروژه‌ها را مستقیم کنترل نمی‌کند.
او فقط به هر پروژه بودجه می‌دهد و نتیجه‌ی آن را در قالب ارزش
کسب‌شده می‌بیند.
فسخ پروژه هم کاری نیست که عامل عمداً انجام دهد، بلکه سازوکاری
درونی در محیط است که با یکی از دو شرط مستقل فعال می‌شود.
شرط اول سررسید قطعی است: اگر دوره‌ی جاری اپیزود به تاریخ پایان
برنامه‌ای به‌علاوه‌ی سقف تأخیر مجاز قرارداد برسد، پروژه بی‌درنگ
فسخ می‌شود.
شرط دوم ضعف عملکرد مداوم است: در هر دوره‌ای که پیشرفت واقعی پروژه
بیش از سقف مجاز قرارداد از پیشرفت برنامه‌ای عقب باشد، یک
شمارنده‌ی تحمل یک واحد کم می‌شود و وقتی این شمارنده به صفر برسد،
پروژه فسخ می‌شود.

فسخ پروژه با پایان اپیزود فرق دارد.
هر پروژه جداگانه فسخ می‌شود و فسخ آن یک جریان نقدی تسویه ایجاد
می‌کند. اپیزود فقط در سه حالت تمام می‌شود: همه‌ی پروژه‌ها یا
تکمیل شده باشند یا فسخ شده باشند، افق کلی برسد، یا بودجه‌ی سبد
تمام شود و هیچ وجهی هم در راه نباشد.

\begin{figure}[htbp]
  \centering
  \includegraphics[width=\textwidth]{figures/feedback_loop.pdf}
  \caption[حلقه‌ی بازخورد یادگیری تقویتی]{راهبری سبد پروژه‌ها به‌صورت
    یک حلقه‌ی بازخورد در یادگیری تقویتی.}
  \label{fig:feedback_loop}
\end{figure}

\subsection{لایه‌های سه‌گانه‌ی مدیریت پروژه}
\label{subsec:three_layers}

برای درک جایگاه کار حاضر در منظومه‌ی مدیریت پروژه، لازم است سه لایه‌ی
اصلی تصمیم‌گیری در این حوزه را با جزئیات بیشتری بررسی کنیم.

\textbf{لایه‌ی استراتژیک} بالاترین سطح تصمیم‌گیری است. در این لایه،
سازمان تصمیم می‌گیرد که اصلاً در کدام حوزه‌های کاری سرمایه‌گذاری
کند، چه ترکیبی از پروژه‌ها با اهداف راهبردی بلندمدتش همخوانی دارد،
و چقدر ریسک کلی می‌تواند بپذیرد. این تصمیم‌ها معمولاً در افق
پنج تا ده ساله اتخاذ می‌شوند و محصول آن‌ها یک سبد اولیه از پروژه‌های
تأیید‌شده است. مدیریت ارشد، هیئت مدیره، و واحدهای برنامه‌ریزی
راهبردی در این لایه نقش دارند.

\textbf{لایه‌ی تاکتیکی} یا همان لایه‌ی مدیریت سبد، جایی است که این
پژوهش در آن قرار می‌گیرد. در این لایه، سبد از پیش انتخاب‌شده باید
اداره شود. تصمیم‌های این لایه شامل تخصیص بودجه‌ی دوره‌ای، اولویت‌بندی
پروژه‌ها در شرایط کمبود منابع، پایش عملکرد پروژه‌ها در برابر
برنامه، و مداخله در صورت بروز انحراف است. افق این تصمیم‌ها
معمولاً یک تا پنج سال است و با چرخه‌ی اجرایی پروژه‌ها همخوانی دارد.
مدیر سبد و تیم دفتر مدیریت پروژه در این لایه کار می‌کنند.

\textbf{لایه‌ی عملیاتی} پایین‌ترین سطح است. در این لایه، هر پروژه
به‌صورت مستقل اداره می‌شود. مدیر پروژه مسئول برنامه‌ریزی فعالیت‌ها،
مدیریت تیم اجرایی، کنترل هزینه در سطح فعالیت، و هماهنگی با
پیمانکاران است. تصمیم‌های این لایه معمولاً افق روزانه تا ماهانه
دارند و بر جزئیات اجرایی تمرکز دارند.

نکته‌ی کلیدی این است که این سه لایه باید با هم هماهنگ باشند، اما
این هماهنگی اغلب دشوار است. مدیر سبد در لایه‌ی تاکتیکی نمی‌تواند
تصمیم‌های راهبردی لایه‌ی بالا را زیر سؤال ببرد، و نمی‌تواند
جزئیات اجرایی لایه‌ی پایین را مستقیماً کنترل کند. او فقط یک اهرم
دارد: تخصیص بودجه. این اهرم قدرتمند است اما به تنهایی کافی نیست
و باید با هوشمندی به کار گرفته شود.

\subsection{جایگاه تخصیص بودجه در لایه‌ی تاکتیکی}
\label{subsec:budget_allocation_role}

در لایه‌ی تاکتیکی، مدیر سبد با مجموعه‌ای از ابزارها کار می‌کند.
پایش عملکرد از طریق شاخص‌های \lr{EVM}، جلسات بازنگری دوره‌ای،
گزارش‌دهی به ذی‌نفعان، و مدیریت ریسک بخشی از این ابزارها هستند.
اما مهم‌ترین اهرم عملیاتی که مدیر سبد در اختیار دارد، تصمیم
تخصیص بودجه است.

تصمیم تخصیص بودجه به دلایل زیر مهم‌ترین ابزار مدیریت سبد است:

اول، \textbf{اثر مستقیم بر پیشرفت پروژه} دارد. در پروژه‌های
عمرانی، بدون پول کار متوقف می‌شود. پیمانکار نمی‌تواند مصالح بخرد،
کارگر نگه دارد، یا ماشین‌آلات اجاره کند اگر پول کافی نداشته باشد.
بنابراین تصمیم چقدر بودجه دادن، مستقیماً تعیین می‌کند که پروژه
چه سرعتی داشته باشد.

دوم، \textbf{اثر غیرمستقیم از طریق ساختار قرارداد} دارد. چون
پرداخت‌ها به نقاط عطف وابسته‌اند، سرعت پیشرفت پروژه تعیین می‌کند
که پیمانکار چه زمانی پول دریافت می‌کند. این پول دوباره به جریان
نقدی سبد برمی‌گردد (از طریق حساب‌های متقابل). بنابراین تصمیم
تخصیص نه‌تنها پول خرج می‌کند بلکه از طریق ساختار قرارداد، زمان
بازگشت پول را هم تعیین می‌کند.

سوم، \textbf{اثر بر ریسک فسخ} دارد. اگر بودجه‌ی کافی داده نشود،
پروژه عقب می‌افتد. اگر این عقب‌افتادگی از حد مجاز قرارداد بگذرد،
شمارنده‌ی تحمل شروع به کم شدن می‌کند. وقتی شمارنده به صفر برسد،
پروژه فسخ می‌شود و هزینه‌های فسخ بر سبد تحمیل می‌شود. بنابراین
تخصیص ناکافی می‌تواند زنجیره‌ای از پیامدهای منفی را آغاز کند.

چهارم، \textbf{اثر بر توزیع منابع در سبد} دارد. هر ریال بودجه‌ای
که به یک پروژه می‌رود از دسترس پروژه‌ی دیگر خارج می‌شود. اگر یک
پروژه به‌طور مداوم بیش از سهم متناسبش از بودجه دریافت کند، پروژه‌های
دیگر ممکن است زیر فشار قرار گیرند. مدیر سبد باید این اثرات متقابل
را در نظر بگیرد.

\subsection{چرا تصمیم‌گیری مرکزی ضروری است}
\label{subsec:centralized_decision}

یک سؤال طبیعی این است: چرا نمی‌توان تخصیص بودجه را به خود مدیران
پروژه واگذار کرد و اجازه داد هر پروژه سهم خود را مطالبه کند؟
پاسخ در ماهیت مشترک بودن منبع نهفته است.

اگر مدیران پروژه تصمیم‌گیران مستقل باشند، هر کدام انگیزه دارند
تا بیشترین سهم ممکن را برای پروژه‌ی خود بگیرند. این رقابت می‌تواند
به چند پیامد ناخواسته منجر شود. اول، پروژه‌هایی که مدیران
قوی‌تر یا پرسروصداتری دارند بودجه‌ی بیشتری می‌گیرند، نه لزوماً
پروژه‌هایی که از نظر اقتصادی بیشتر نیاز دارند. دوم، مدیران
پروژه انگیزه دارند تا وضعیت پروژه‌ی خود را بهتر از واقع نشان
دهند تا بودجه‌ی بیشتری دریافت کنند، که این اطلاعات نامتقارن
را تشدید می‌کند. سوم، کسی از منظر کل سبد تصمیم نمی‌گیرد و
ممکن است تخصیص بودجه به‌شکلی باشد که برای هر پروژه‌ی منفرد
معقول اما برای کل سبد زیان‌آور باشد.

تصمیم‌گیری مرکزی این مشکلات را حل می‌کند اما مشکل جدیدی ایجاد
می‌کند: مدیر سبد باید اطلاعات کافی و ابزار مناسب برای تصمیم‌گیری
بهینه داشته باشد. این پژوهش بخشی از این ابزار را فراهم می‌کند.

% ------------------------------------------------------------
\section{سازوکارهای قرارداد و ساختار جریان نقدی}
\label{sec:contract_mechanics}
% ------------------------------------------------------------

رابطه‌ی مالی مدیر سبد با هر پیمانکار طبق فرم‌های قرارداد
\lr{FIDIC} \citep{FIDIC2017} تنظیم می‌شود.
سه سازوکار ساختار جریان نقدی را مشخص می‌کنند.

\paragraph{پرداخت بر پایه‌ی نقطه‌ی عطف.}
پرداخت به رسیدن پروژه به آستانه‌های پیشرفتِ تأییدشده بستگی دارد.
پرداخت ناخالص در نقطه‌ی عطف $j$ کسر ثابتی از قیمت قرارداد است که
در خود قرارداد تعیین شده است.
چون پرداخت به پیشرفتِ تأییدشده وابسته است، زمان ورود پول به
عملکرد واقعی پیمانکار و به تصمیم‌های تخصیص مدیر سبد بستگی دارد.
در مدل ما عدم‌قطعیت جریان نقدی فقط از یک منبع می‌آید:
بازده تصادفی بودجه به پیشرفت، یعنی $\eta_{i,t}$.
این بازده تعیین می‌کند که یک مقدار مشخص بودجه در هر دوره
چقدر پیشرفت فیزیکی ایجاد کند.
بهره‌وریِ نامطمئن، زمان رسیدن به نقاط عطف را نامطمئن می‌کند و
این هم ورود پول را نامطمئن می‌کند. به همین دلیل نیازی به مدل جداگانه
برای تأخیر در پرداخت نیست.

\paragraph{پیش‌پرداخت و بازپس‌گیری تدریجی.}
در شروع پروژه، پیش‌پرداختی برابر کسر ثابتی از قیمت قرارداد به
پیمانکار داده می‌شود.
این مبلغ با کسر درصدی از صورت‌وضعیت‌های بعدی به‌تدریج بازپس
گرفته می‌شود. آستانه‌ی شروع کسر و نرخ آن، پارامترهای قراردادیِ
هر پروژه هستند \citep{FIDIC2017}.

\paragraph{سپرده‌ی حسن انجام کار.}
کسری از هر پرداخت نقطه‌ی عطف به‌عنوان تضمین عملکرد نگه داشته
می‌شود و پس از تکمیل پروژه یک‌جا آزاد می‌شود.
به این ترتیب خود ساختار پرداخت، پیمانکار را به تکمیل پروژه تشویق
می‌کند.

\paragraph{مدیریت ارزش کسب‌شده به‌عنوان سیگنال عملکرد.}
چارچوب مدیریت ارزش کسب‌شده (\lr{EVM}) \citep{Fleming2010} زبان
لازم برای پایش عملکرد پیمانکار را فراهم می‌کند.
شاخص عملکرد هزینه، کارایی هزینه را می‌سنجد و شاخص عملکرد
زمان‌بندی بر پایه‌ی زمان کسب‌شده، کارایی زمانی را.
برآورد هزینه در زمان تکمیل با مخرجی که هم هزینه و هم زمان‌بندی را
در نظر می‌گیرد، کل هزینه‌ی پروژه را پیش‌بینی می‌کند.
شاخص عملکرد لازم برای تکمیل نیز نشان می‌دهد برای تمام کردن کار
باقی‌مانده در چارچوب بودجه، چه کارایی‌ای لازم است.
این چهار شاخص درون محیط محاسبه می‌شوند و در هر دوره در رکورد
حسابرسی ثبت می‌شوند، اما مستقیماً در بردار مشاهده قرار ندارند.
نقش آن‌ها واسطه است: ویژگی‌های مربوط به تخصیص جبرانی و نقطه‌ی عطف
هدف از روی آن‌ها ساخته می‌شوند.

\paragraph{فسخ قرارداد.}
ویژگی مهم محیط پیشنهادی این است که فسخ پروژه نتیجه‌ی خودِ
تصمیم‌های مالی مدیر است و مثل یک شوک بیرونی به سبد وارد نمی‌شود.
در عمل، پیمانکار می‌تواند پروژه را از دو راه کاملاً متفاوت از
دست بدهد.
راه اول تدریجی است: کم‌بودجه‌بودن مداوم باعث می‌شود تأخیر از حد
مجاز قرارداد بگذرد و دوره‌ی مهلتی هم که برای جبران در نظر گرفته
شده تمام شود.
راه دوم ناگهانی است: تاریخ پایان پروژه از حد مجاز قرارداد بگذرد،
حتی اگر پروژه تا جبران عقب‌ماندگی فاصله‌ی کمی داشته باشد.
مدل‌های تک‌پروژه‌ای قبلی معمولاً فقط یکی از این دو حالت را در نظر
می‌گیرند و به همین دلیل ریسکی را که مدیر واقعاً با آن روبه‌روست
کمتر از مقدار واقعی نشان می‌دهند.
فسخ پیامد مالی مستقیم هم دارد که باید با وضعیت نقدی سبد تسویه
شود، چون کاری که پولش پرداخت شده و کاری که واقعاً انجام شده،
در لحظه‌ی فسخ تقریباً هیچ‌وقت دقیقاً برابر نیستند.
در محیط ما هر دو شرط نقض، مهلت مربوط به آن‌ها و جریان نقدی تسویه
مستقیماً در دینامیک گذار قرار داده شده‌اند.
بنابراین هر سیاستی ناچار است خطر فسخ را نتیجه‌ی انتخاب‌های
تخصیصی خودش بداند، که دقیقاً همان وضعیتی است که یک مدیر سبد
واقعی با آن مواجه است.

\subsection{اهمیت صحیح مدل کردن ساختار قرارداد}
\label{subsec:contract_modeling_importance}

یکی از نوآوری‌های اساسی این پژوهش نسبت به کارهای قبلی، توجه دقیق
به جزئیات ساختار قراردادی است. بسیاری از مدل‌های موجود ساختار
پرداخت را ساده می‌کنند: فرض می‌کنند که بودجه‌ی تخصیص‌یافته
مستقیماً به پیشرفت تبدیل می‌شود و پرداخت‌های برگشتی وجود ندارند.
این ساده‌سازی در مواردی که هدف فقط مدل کردن پیشرفت فیزیکی است
قابل قبول است، اما وقتی هدف بهینه‌سازی جریان نقدی است، این
ساده‌سازی اثرات مهمی را از دست می‌دهد.

پیش‌پرداخت قراردادی یکی از این اثرات است. وقتی پروژه‌ای شروع
می‌شود، مدیر سبد مبلغ قابل توجهی به‌عنوان پیش‌پرداخت خارج می‌کند.
این پول در ترازنامه به‌عنوان طلب از پیمانکار ثبت می‌شود و با
پیشرفت پروژه به‌تدریج کسر می‌شود. نادیده گرفتن این سازوکار
به این معناست که جریان نقدی خروجی اولیه‌ی پروژه را دست کم بگیریم
و زمان بازیابی آن را نادیده بگیریم.

سپرده‌ی حسن انجام کار اثر دیگری دارد. این کسر از هر پرداخت نقطه‌ی
عطف، به این معناست که پیمانکار در طول اجرا کمتر از آنچه می‌خواند
دریافت می‌کند. این کمبود نقدینگی پیمانکار ممکن است عملکردش را
تحت‌تأثیر قرار دهد، به‌خصوص در مراحل پایانی پروژه که کارهای باقی‌مانده
بیشتر از هزینه‌ی پرداخت‌نشده‌اند. مدل باید این پیامدهای
رفتاری را هم لحاظ کند.

در نهایت، ساختار فسخ قرارداد مهم‌ترین اثر را دارد. وقتی پروژه‌ای
فسخ می‌شود، تسویه‌ی مالی آن می‌تواند به نفع یا ضرر مدیر سبد
باشد، بسته به اینکه مجموع پرداخت‌های انجام‌شده چقدر با کار
واقعاً انجام‌شده تفاوت دارد. اگر مدیر سبد قبل از فسخ پول
زیادی داده باشد اما پیشرفت کمی حاصل شده باشد، تسویه ممکن است
جریان نقدی مثبتی برای سبد ایجاد کند. اگر برعکس باشد، تسویه
جریان نقدی منفی اضافی تحمیل می‌کند. این پیامدهای مالی فسخ
باید دقیقاً در مدل لحاظ شوند.

\subsection{نقش \lr{EVM} در راهبری سبد}
\label{subsec:evm_in_portfolio}

مدیریت ارزش کسب‌شده در اصل برای کنترل پروژه‌های منفرد توسعه
یافته، اما در سطح سبد نیز کاربرد مهمی دارد. در سطح سبد، شاخص‌های
\lr{EVM} اطلاعاتی فراهم می‌کنند که مدیر سبد می‌تواند از آن‌ها
برای تشخیص پروژه‌های مشکل‌دار استفاده کند.

شاخص عملکرد هزینه که کمتر از یک باشد نشان می‌دهد پروژه برای هر
واحد کار انجام‌شده بیش از بودجه هزینه می‌کند. اگر این شاخص در
چند دوره‌ی متوالی پایین‌تر از حد قابل قبول قرارداد باقی بماند،
این یک هشدار جدی است. به همین ترتیب، شاخص عملکرد زمان‌بندی که
از یک پایین‌تر باشد نشان می‌دهد پروژه کندتر از برنامه پیش
می‌رود.

ترکیب این دو شاخص یک تصویر جامع از وضعیت پروژه می‌دهد. پروژه‌ای
که هم از نظر هزینه و هم از نظر زمان‌بندی ضعیف است، مدیر سبد را
با یک تصمیم سخت روبه‌رو می‌کند: آیا باید با تخصیص بودجه‌ی بیشتر
سعی در نجات آن کرد، یا این کار فقط منابع بیشتری را در یک پروژه‌ی
محکوم‌به‌شکست هدر می‌دهد؟

پاسخ به این سؤال بستگی به وضعیت کل سبد دارد. اگر پروژه‌های دیگر
سبد در شرایط خوبی هستند و بودجه‌ی کافی موجود است، شاید ارزش داشته
باشد که پول بیشتری برای نجات پروژه‌ی مشکل‌دار صرف شود. اما اگر
بودجه‌ی سبد تنگ است، کمک به یک پروژه ممکن است بقیه را به خطر بیندازد.
این تبادل بین پروژه‌ها است که مدیریت سبد را از مدیریت پروژه‌ی منفرد
متمایز می‌کند، و همین تبادل است که مدل‌سازی صحیح آن را ضروری می‌سازد.

% ------------------------------------------------------------
\section{شکاف بهینگی و دلیل نیاز به سیاست‌های یادگیری‌شده}
\label{sec:gap}
% ------------------------------------------------------------

یک پرسش تجربی اساسی این است که آیا مسئله آن‌قدر دشوار هست که
استفاده از سیاست‌های یادگیری‌شده را توجیه کند.
برای پاسخ، سه روش ابتکاری حریصانه‌ی قطعی را با یک جواب بهینه‌ی
اطلاعات‌کامل، که به‌صورت دستی محاسبه شده، مقایسه می‌کنیم.
این مقایسه روی چهار پیکربندی آزمون سبد انجام می‌شود که پیچیدگی
آن‌ها پله‌پله بیشتر می‌شود.

روش \textbf{رسیدن به برنامه} (\lr{Reach-Plan}، $\pi^{\mathrm{rp}}$)
در هر دوره دقیقاً آن‌قدر بودجه می‌دهد که پیشرفت هر پروژه‌ی فعال به
پیشرفت برنامه‌ای روی منحنی \lr{S} برسد و فاصله‌ی پیشرفت کاملاً بسته
شود.
روش \textbf{جبران \lr{T}} (\lr{Catchup-T}، $\pi^{\mathrm{cu}}$)
فقط همان‌قدر بودجه می‌دهد که تأخیر پیشرفت در حد مجاز قرارداد بماند،
و به این ترتیب هزینه‌ها را کمینه می‌کند.
روش \textbf{جبران \lr{T} بعدی} (\lr{Catchup-Next-T}،
$\pi^{\mathrm{cu}+1}$) محتاط‌تر است: همین امروز آن‌قدر بودجه
می‌دهد که تأخیر پیش‌بینی‌شده برای دوره‌ی بعد هم در حد مجاز بماند.
\textbf{جواب بهینه‌ی اطلاعات‌کامل} ($\pi^*$) به‌صورت دستی و با فرض
اینکه همه‌ی مقدارهای تصادفی از قبل معلوم‌اند محاسبه می‌شود و
نقش کران بالا را دارد.

شکل~\ref{fig:gap_plot} ارزش خالص فعلی (\lr{NPV}) تجمعی اپیزود را
برای هر روش در چهار پیکربندی نشان می‌دهد.
در ساده‌ترین پیکربندی (یک پروژه با بازده ثابت) روش‌های حریصانه با
اختلاف کمی به جواب بهینه نزدیک می‌شوند.
هرچه سبد پیچیده‌تر می‌شود، یعنی پروژه‌ها بیشتر می‌شوند، شروع آن‌ها
در زمان‌های مختلف قرار می‌گیرد، جریان‌های نقدی بر هم اثر می‌گذارند
و قیدهای ترتیب نقاط عطف اضافه می‌شود، این شکاف به‌طور منظم بزرگ‌تر
می‌شود.
در پیچیده‌ترین پیکربندی، بهترین روش حریصانه بیش از ۶۰ واحد
\lr{NPV} از جواب بهینه عقب است و یکی از روش‌ها حتی به نتیجه‌ی
زیان‌ده می‌رسد.
این افزایش یکنواخت نشان می‌دهد که مسئله‌ی تخصیص در مقیاس بزرگ
واقعاً دشوار می‌شود و برای عامل‌های یادگیری تقویتی هدف مشخصی برای
بهبود وجود دارد.

\subsection{چرا روش‌های حریصانه کافی نیستند}
\label{subsec:why_greedy_fails}

روش‌های حریصانه‌ی قطعی از نظر شهودی جذاب هستند. منطق آن‌ها ساده
و قابل توضیح است: یا برنامه را دنبال کن، یا از فسخ جلوگیری کن،
یا یک قدم جلوتر از مشکل باش. این سادگی یک مزیت عملی واقعی است؛
مدیران می‌توانند این قوانین را بدون محاسبات پیچیده اعمال کنند.

اما روش‌های حریصانه در مواجهه با چند چالش اساسی ناکام می‌مانند.

\textbf{چالش اول: کوته‌بینی زمانی.} روش حریصانه بر اساس وضعیت
فعلی تصمیم می‌گیرد، نه بر اساس پیامدهای آینده. مثلاً روش
«رسیدن به برنامه» در هر دوره پول کافی می‌دهد تا عقب‌ماندگی پر
شود. اما این ممکن است در شرایطی که بودجه‌ی کل محدود است، پول
را به پروژه‌ای بدهد که نقطه‌ی عطف بعدی‌اش هنوز دور است، در حالی
که پروژه‌ی دیگری نزدیک به نقطه‌ی عطف است و با کمی پول بیشتر
می‌توانست آن را تریگر کند و پول قابل توجهی برگرداند.

\textbf{چالش دوم: نادیده گرفتن تعاملات بین پروژه‌ای.} در یک سبد
چندپروژه‌ای، تصمیم برای یک پروژه بر دیگران اثر می‌گذارد. روش‌های
حریصانه‌ی ساده هر پروژه را مستقل در نظر می‌گیرند و این تعاملات
را نادیده می‌گیرند. اما در واقعیت، بودجه‌ی اضافه‌ای که به
پروژه‌ی الف داده می‌شود از دهان پروژه‌ی ب می‌رود، و ممکن است
این انتقال در مجموع به سبد آسیب بزند.

\textbf{چالش سوم: ناتوانی در تطبیق با تحقق‌های تصادفی.} بازده‌ی
بودجه به پیشرفت در هر دوره متفاوت است. روش‌های حریصانه قانون
ثابتی دارند که با این تغییرپذیری تطبیق نمی‌دهند. یک سیاست یادگرفته‌شده
می‌تواند یاد بگیرد که در دوره‌هایی که بازده تاریخاً بالا بوده،
تخصیص بیشتری انجام دهد تا از فرصت نقطه‌ی عطف استفاده کند، و
در دوره‌هایی با بازده پایین، محتاط‌تر باشد.

\textbf{چالش چهارم: بدون آستانه‌ی تصمیم فسخ.} روش‌های حریصانه
ابزاری برای تشخیص اینکه «ادامه دادن به این پروژه دیگر صرف
نمی‌کند» ندارند. آن‌ها تا آخر به قوانین ثابت خود پایبند هستند
حتی وقتی واضح است که پروژه‌ای محکوم به شکست است و بودجه‌ی آن
بهتر است به جای دیگری برود. یک سیاست یادگرفته‌شده می‌تواند یاد
بگیرد که در چنین شرایطی تخصیص را کاهش دهد و بگذارد فسخ با
شرایط مطلوب‌تری اتفاق بیفتد.

\subsection{اهمیت کران بالا در ارزیابی سیاست}
\label{subsec:upper_bound_importance}

یکی از نوآوری‌های روش‌شناختی این پژوهش، محاسبه‌ی جواب بهینه‌ی
اطلاعات‌کامل به‌عنوان کران بالاست. این کران بالا نشان می‌دهد که
بهترین سیاست ممکن، حتی اگر آینده را کاملاً بداند، چه عملکردی
می‌تواند داشته باشد.

فاصله‌ی یک سیاست از این کران بالا معیاری دقیق برای کیفیت آن
سیاست است. روش‌های حریصانه از این کران فاصله‌ی قابل اندازه‌گیری
دارند. سیاست‌های یادگرفته‌شده باید بتوانند این فاصله را کاهش دهند.
اما به دست آوردن این کران بالا در مسائل واقعی دشوار است، چون
حتی دانستن آینده هم مسئله‌ی بهینه‌سازی پیچیده‌ای ایجاد می‌کند.

در مدل ما، چون ساختار قرارداد و پارامترهای تصادفی کاملاً مشخص است،
می‌توان با داشتن تحقق‌های تصادفی از پیش مشخص، برنامه‌ریزی پویای
دقیقی انجام داد و بهینه‌ی اطلاعات‌کامل را محاسبه کرد. این
امکان در مسائل واقعی وجود ندارد، اما در محیط شبیه‌سازی ما
قابل دستیابی است و به ما اجازه می‌دهد که شکاف بهینگی را دقیقاً
اندازه‌گیری کنیم.

% ------------------------------------------------------------
\section{چارچوب یادگیری تقویتی برای مسئله‌ی تخصیص}
\label{sec:rl_framework}
% ------------------------------------------------------------

یادگیری تقویتی چارچوبی است که در آن یک عامل یاد می‌گیرد چگونه
در یک محیط رفتار کند تا یک هدف عددی را بیشینه کند. این چارچوب
سه عنصر اصلی دارد: عامل که تصمیم می‌گیرد، محیط که به تصمیم‌ها
واکنش نشان می‌دهد، و سیگنال پاداش که نشان می‌دهد هر تصمیم
چقدر خوب بود.

در مسئله‌ی تخصیص بودجه‌ی سبد پروژه، این سه عنصر به‌طور طبیعی
تعریف می‌شوند. عامل همان مدیر سبد است که در هر دوره تخصیص
بودجه را تعیین می‌کند. محیط شامل تمام پروژه‌های فعال در سبد
است که با پارامترهای قراردادی و بازده‌های تصادفی مشخص می‌شوند.
سیگنال پاداش جریان نقدی خالص تنزیل‌شده‌ی دریافتی در هر دوره است.

\subsection{فرایند تصمیم مارکوف به‌عنوان صورت‌بندی رسمی}
\label{subsec:mdp_formulation_intro}

فرایند تصمیم مارکوف یا \lr{MDP} صورت‌بندی ریاضی استاندارد برای
مسائل تصمیم‌گیری ترتیبی است. یک \lr{MDP} از چهار جزء تشکیل می‌شود:
فضای حالت که وضعیت محیط را توصیف می‌کند، فضای عمل که مجموعه
تصمیم‌های ممکن است، تابع گذار که نشان می‌دهد هر تصمیم در هر حالت
به چه حالت بعدی منجر می‌شود، و تابع پاداش که ارزش هر گذار را مشخص
می‌کند.

صورت‌بندی مسئله‌ی تخصیص بودجه‌ی سبد به‌صورت یک \lr{MDP} چند مزیت
مهم دارد. اول، این صورت‌بندی وابستگی زمانی تصمیم‌ها را به‌صورت
صریح مدل می‌کند؛ حالت در هر دوره شامل اطلاعات کافی برای تصمیم‌گیری
بهینه است. دوم، عدم‌قطعیت از طریق احتمال گذار به‌طور طبیعی وارد
مدل می‌شود. سوم، هدف بیشینه‌کردن مجموع تنزیل‌شده‌ی پاداش‌ها
دقیقاً همان هدفی است که مدیر سبد دنبال می‌کند.

صورت‌بندی \lr{MDP} همچنین اجازه می‌دهد از طیف گسترده‌ای از
الگوریتم‌های یادگیری تقویتی استفاده شود. این الگوریتم‌ها از
تعامل با محیط یاد می‌گیرند و به‌مرور سیاستی پیدا می‌کنند که
پاداش کل را بیشینه کند. آموزش این عامل‌ها روی محیط شبیه‌سازی‌شده‌ی
این پژوهش گام طبیعی بعدی است.

\subsection{چرا یادگیری تقویتی و نه برنامه‌ریزی تصادفی کلاسیک}
\label{subsec:rl_vs_sp}

شاید بپرسیم که اگر مسئله را می‌توان به‌صورت \lr{MDP} صورت‌بندی کرد،
چرا از روش‌های حل دقیق مثل برنامه‌ریزی پویا استفاده نمی‌شود؟

پاسخ در ابعاد مسئله نهفته است. فضای حالت مسئله‌ی تخصیص بودجه‌ی
سبد با تعداد پروژه‌ها به‌صورت نمایی رشد می‌کند. برای یک سبد با
پنج پروژه، هر پروژه چندین متغیر حالت دارد که شامل پیشرفت فعلی،
فاصله از نقطه‌ی عطف بعدی، شمارنده‌ی مهلت جبران، موجودی پیش‌پرداخت
و چند متغیر دیگر است. ترکیب این متغیرها برای پنج پروژه و موجودی
نقدی سبد، یک فضای حالت چندبعدی بزرگ می‌سازد که برنامه‌ریزی پویای
دقیق برای آن محاسباتاً غیرعملی است.

برنامه‌ریزی تصادفی چندمرحله‌ای با رویکرد درخت سناریو نیز با
مشکل مشابهی روبه‌روست. تعداد سناریوهای ممکن برای یک افق برنامه‌ریزی
طولانی با چند پروژه به‌طور نمایی رشد می‌کند و این رویکرد را برای
مسائل بزرگ غیرعملی می‌کند.

یادگیری تقویتی از طریق تعمیم بین حالت‌های مشابه، با این لعنت
ابعاد کنار می‌آید. یک عامل یادگیری تقویتی یاد می‌گیرد که در
حالت‌های مشابه رفتارهای مشابهی داشته باشد، بدون اینکه نیاز باشد
هر حالت ممکن را جداگانه ببیند. این توانایی تعمیم چیزی است که
روش‌های دقیق ندارند و دقیقاً همان چیزی است که برای مسائل بزرگ
واقعی نیاز داریم.

% ------------------------------------------------------------
\section{شکاف بهینگی و دلیل نیاز به سیاست‌های یادگیری‌شده --- ادامه}
\label{sec:gap_continued}
% ------------------------------------------------------------

\subsection{تفسیر کمّی شکاف}
\label{subsec:gap_interpretation}

نتیجه‌ی اصلی آزمون خط مبنا این است که شکاف بهینگی با پیچیدگی سبد
به‌طور یکنواخت بزرگ‌تر می‌شود. این نتیجه صرفاً یک مشاهده‌ی
تجربی نیست؛ از ماهیت مسئله هم قابل استدلال است.

وقتی فقط یک پروژه وجود دارد، مسئله ساده است: هر چه بودجه‌ی بیشتری
بدهی پیشرفت بیشتری می‌گیری. روش‌های حریصانه در این حالت تقریباً
بهینه‌اند چون فضایی برای بهبود وجود ندارد.

وقتی دو پروژه وجود داشته باشند، تبادل بین آن‌ها ایجاد می‌شود. اکنون
باید تصمیم گرفت که کدام پروژه در این دوره بودجه‌ی بیشتری بگیرد. این
تصمیم بستگی دارد به اینکه کدام پروژه به نقطه‌ی عطف نزدیک‌تر است،
کدام در خطر فسخ است، و کدام ارزش افزوده‌ی بالاتری دارد. روش‌های
حریصانه قادر به در نظر گرفتن همه‌ی این عوامل به‌طور هم‌زمان نیستند.

با افزایش تعداد پروژه‌ها به سه، چهار، پنج، تعداد تصمیم‌های ممکن
در هر دوره به‌طور نمایی رشد می‌کند و پیچیدگی تعاملات بین پروژه‌ها
هم بیشتر می‌شود. روش‌های حریصانه در این پیچیدگی گم می‌شوند و
ضرری که از تصمیم‌های سرانگشتی ایجاد می‌شود با تعداد پروژه‌ها
بزرگ‌تر می‌شود.

این تحلیل نشان می‌دهد که نیاز به روش‌های هوشمندتر در سبدهای
بزرگ‌تر بیشتر است. شرکت‌هایی که ده‌ها پروژه را هم‌زمان مدیریت
می‌کنند، پتانسیل بیشتری برای بهبود از طریق تخصیص بهینه دارند
تا شرکت‌هایی با دو یا سه پروژه.

\subsection{مفهوم سیاست یادگرفته‌شده}
\label{subsec:learned_policy_concept}

یک سیاست در ادبیات یادگیری تقویتی تابعی است که از حالت محیط به
عمل نگاشت می‌کند. یک سیاست یادگرفته‌شده سیاستی است که از طریق
تعامل با محیط و دریافت پاداش‌ها به‌دست آمده، نه از طریق طراحی
قاعده‌ی صریح.

مزیت اصلی سیاست‌های یادگرفته‌شده نسبت به قوانین سرانگشتی این
است که این سیاست‌ها می‌توانند ساختارهای پیچیده‌ای را یاد بگیرند
که بیان صریح آن‌ها دشوار است. یک مدیر سبد باتجربه ممکن است
بداند که «وقتی پروژه‌ی الف نزدیک نقطه‌ی عطف است و پروژه‌ی ب
هنوز دور است، بهتر است بودجه‌ی بیشتری به الف بدهم». اما بیان
این قانون به‌صورت صریح برای همه‌ی حالات ممکن، دشوار است. یک
عامل یادگیری تقویتی می‌تواند این الگو را از تجربه یاد بگیرد
بدون اینکه کسی آن را صریحاً برنامه‌ریزی کرده باشد.

محیط شبیه‌سازی‌شده‌ای که این پژوهش ارائه می‌دهد، دقیقاً همان بستری
است که چنین سیاست‌هایی می‌توانند در آن آموزش ببینند. هر بار که
عامل در محیط اجرا می‌شود، اطلاعات جدیدی درباره‌ی اینکه کدام
تخصیص‌ها نتایج بهتری دارند به دست می‌آورد. با هزاران اپیزود
آموزش، این اطلاعات می‌توانند به یک سیاست تبدیل شوند که به‌طور
قابل اندازه‌گیری از روش‌های حریصانه بهتر عمل می‌کند.

% ------------------------------------------------------------
\section{دستاوردها و ساختار پایان‌نامه}
\label{sec:organisation}
% ------------------------------------------------------------

این پایان‌نامه سه دستاورد اصلی دارد.

\begin{enumerate}[leftmargin=*, label=\textbf{د\arabic*.}]

    \item \textbf{محیط.}
    یک محیط \lr{MDP} قابل پیکربندی و سازگار با \lr{Gymnasium} که
    چرخه‌ی کامل جریان نقدی قرارداد \lr{FIDIC} را در خود دارد:
    پیش‌پرداخت و بازپس‌گیری تدریجی آن، پرداخت‌های نقطه‌ی عطف که
    بر پایه‌ی پیشرفت وزن‌دهی می‌شوند، کسر و آزادسازی سپرده‌ی حسن
    انجام کار، بهره‌ی برداشت از خزانه برای موجودی نقدی منفی پروژه،
    و یک دروازه‌ی فسخ با دو شرط که برای هر پروژه شمارنده‌ی مهلت
    جبران جداگانه‌ای دارد.
    همه‌ی پارامترهای محیط از راه یک دیکشنری پیکربندی تعیین
    می‌شوند که مقدارهای آن توزیع‌های احتمال هستند.
    به این ترتیب فعالان هر حوزه‌ی پروژه‌محور، مانند ساخت‌وساز،
    فناوری اطلاعات یا پژوهش و توسعه، می‌توانند با برازش این توزیع‌ها
    روی داده‌های تاریخی شرکت خود، محیطی مخصوص حوزه‌ی خود بسازند.
    سازگاری با \lr{API} محیط \lr{Gymnasium} با ۵۱ بررسی در ۱۰ بخش
    آزمون تأیید شده است و گواهی کامل آن در
    پیوست~\ref{app:gym_certificate} آمده است.
    محیط یک دفتر ثبت حسابرسی \lr{SQLite} هم دارد که در هر گام
    اپیزود، هر تخصیص، هر رویداد جریان نقدی و هر پرچم نقض را ثبت
    می‌کند. به این ترتیب رفتار سیاست کاملاً قابل بررسی و تفسیر
    است.

    \item \textbf{صورت‌بندی \lr{MDP}.}
    تعریف رسمی مسئله‌ی بودجه‌بندی سبد به‌صورت یک فرایند تصمیم‌گیری
    مارکوف، شامل فضای حالت، فضای عمل، دینامیک گذار، سیگنال پاداش و
    بردار مشاهده در قالب فرمول‌های بسته.
    مشخصات ریاضی و محیط پیاده‌سازی‌شده کاملاً با هم هماهنگ‌اند.

    \item \textbf{آزمون خط مبنا و شکاف بهینگی.}
    ارزیابی سه روش ابتکاری حریصانه‌ی قطعی در برابر جواب بهینه‌ی
    اطلاعات‌کامل روی چهار پیکربندی آزمون.
    نتیجه‌ی این ارزیابی یک شکاف بهینگی کمّی است که با پیچیدگی سبد
    به‌طور یکنواخت بزرگ‌تر می‌شود و هدف عملکردی عامل‌های یادگیری
    تقویتی آینده را مشخص می‌کند.

\end{enumerate}

\subsection{اهمیت آزمون‌پذیری و قابلیت تکرار}
\label{subsec:reproducibility}

یکی از ارزش‌های کمتر دیده‌شده‌ی این پژوهش، توجه به آزمون‌پذیری
و قابلیت تکرار است. محیط شبیه‌سازی‌شده نه فقط برای آموزش
عامل‌های یادگیری تقویتی، بلکه به‌عنوان یک بنچمارک استاندارد
برای مقایسه‌ی روش‌های مختلف طراحی شده است.

هر محقق یا مدیری که می‌خواهد روش جدیدی برای تخصیص بودجه‌ی سبد
پیشنهاد دهد، می‌تواند آن را روی همین محیط ارزیابی کند و نتایج را
با روش‌های خط مبنا مقایسه کند. این قابلیت مقایسه‌ی استاندارد
چیزی است که در ادبیات مدیریت پروژه کم‌تر وجود داشته.

گواهی سازگاری با \lr{Gymnasium} که ۵۱ آزمون را پوشش می‌دهد،
تضمین می‌کند که محیط رفتار پیش‌بینی‌پذیری دارد و با ابزارهای
استاندارد یادگیری تقویتی سازگار است. دفتر ثبت حسابرسی \lr{SQLite}
امکان بازبینی دقیق رفتار هر سیاست را فراهم می‌کند و اجازه
می‌دهد که هر نتیجه کاملاً توضیح داده شود.

\subsection{راهنمای خواندن پایان‌نامه}
\label{subsec:reading_guide}

فصل~\ref{chap:literature} ادبیات مرتبط را در چهار محور مرور می‌کند:
مبانی مدیریت پروژه، مدیریت و انتخاب سبد پروژه، بودجه‌بندی و
مدل‌سازی جریان نقدی در مدیریت سبد پروژه، و یادگیری تقویتی برای
تخصیص ترتیبی منابع.
فصل~\ref{chap:problem} صورت‌بندی \lr{MDP} را به‌طور رسمی بیان
می‌کند.
فصل~\ref{chap:evaluation} ارزیابی روی پیکربندی‌های آزمون را ارائه
می‌دهد: سه روش ابتکاری حریصانه‌ی قطعی و یک جواب بهینه‌ی
اطلاعات‌کامل روی چهار پیکربندی سبد، که شکاف بهینگی لازم برای
توجیه عامل‌های یادگیری‌شده را نشان می‌دهد.
آموزش عامل‌های یادگیری تقویتی روی این محیط گام طبیعی بعدی است و
در کارهای آینده به آن پرداخته می‌شود.

خواننده‌ای که با مدیریت پروژه آشناست اما با یادگیری تقویتی آشنایی
کمتری دارد، توصیه می‌شود فصل دوم را از بخش یادگیری تقویتی آن آغاز
کند تا تصویری اولیه از چارچوب داشته باشد، سپس به فصل سوم بازگردد.
خواننده‌ای که از حوزه‌ی یادگیری ماشین می‌آید، توصیه می‌شود بخش
سازوکارهای قرارداد در همین فصل را با دقت بیشتری مطالعه کند، چون
این بخش شرایط مرزی محیط را توضیح می‌دهد و درک درست آن برای
تفسیر نتایج ضروری است.